\DocumentMetadata{
  pdfversion=2.0,
  pdfstandard=ua-2,
  lang=en-US,
  tagging=on,
  tagging-setup={math/setup=mathml-SE,tabsorder=structure}
}
\documentclass[11pt,letterpaper]{article}
\usepackage[margin=0.68in,includefoot]{geometry}
\usepackage{newcomputermodern}
\usepackage{amsmath,amscd,mathtools}
\usepackage{tikz,adjustbox}
\providecommand{\square}{\mdlgwhtsquare}
\usepackage[hidelinks]{hyperref}
\hypersetup{
  pdftitle={Algebra First-Year Exam, May 2022, Part 1},
  pdfauthor={University of Florida Department of Mathematics},
  pdfsubject={Graduate mathematics examination},
  pdfdisplaydoctitle=true,
  bookmarksopen=true
}
\setcounter{secnumdepth}{0}
\setlength{\parindent}{0pt}
\setlength{\parskip}{0.28em}
\setlength{\emergencystretch}{3em}
\setlength{\leftmargini}{2.15em}
\setlength{\leftmarginii}{2.65em}
\setlength{\leftmarginiii}{3em}
\pagestyle{empty}
\raggedbottom
\allowdisplaybreaks
\NewTaggingSocketPlug{block/list/label}{examproblem}{%
  \tagstructbegin{tag=Lbl}\tagstructend%
  \tagstructbegin{tag=\LBody}%
  \tagstructbegin{tag=\ProblemHeadingTag,title-o={\ProblemHeadingText}}%
  \tagmcbegin{tag=\ProblemHeadingTag,actualtext={\ProblemHeadingText}}%
  #2%
  \tagmcend%
  \tagstructend%
}
\newenvironment{examproblems}[2]{%
  \def\ProblemHeadingTag{#2}%
  \AssignTaggingSocketPlug{block/list/label}{examproblem}%
  \begin{enumerate}%
  \setcounter{enumi}{#1}%
  \renewcommand{\labelenumi}{\textbf{\theenumi.}}%
  \setlength{\topsep}{0.35em}%
  \setlength{\partopsep}{0pt}%
  \setlength{\itemsep}{0.48em}%
  \setlength{\parsep}{0.18em}%
}{\end{enumerate}}
\newenvironment{examsubparts}{%
  \AssignTaggingSocketPlug{block/list/label}{default}%
  \begin{enumerate}%
  \setlength{\topsep}{0.18em}%
  \setlength{\partopsep}{0pt}%
  \setlength{\itemsep}{0.16em}%
  \setlength{\parsep}{0.1em}%
}{\end{enumerate}}
\newenvironment{examinstructions}{%
  \par\begingroup\itshape
}{%
  \par\endgroup\vspace{0.18em}
}
\makeatletter
\renewcommand\subsection{\@startsection{subsection}{2}{0pt}%
  {-1.25ex plus -.25ex minus -.1ex}%
  {0.55ex plus .1ex}%
  {\normalfont\large\bfseries}}
\renewcommand{\@maketitle}{%
  \newpage\null\vskip -1em%
  {\footnotesize\bfseries
    \href{https://www.ufl.edu/}{University of Florida}, \href{https://math.ufl.edu/}{Department of Mathematics}\par}%
  \vskip -0.5em%
  \tagstructbegin{tag=H1,title={Algebra First-Year Exam, May 2022, Part 1}}%
  \tagpdfparaOff%
  \tagmcbegin{tag=H1}%
    {\LARGE\bfseries\href{https://gma.math.ufl.edu/exams/algebra-first-year/}{Algebra First-Year Exam}, May 2022, Part 1\par}%
  \tagmcend%
  \tagpdfparaOn%
  \tagstructend%
  \vskip 0.25em%
  \tagmcbegin{artifact}%
    \hrule height 1pt%
  \tagmcend%
  \vskip 1em%
}
\makeatother
\title{Algebra First-Year Exam, May 2022, Part 1}
\author{}
\date{}
\begin{document}
\maketitle
\thispagestyle{empty}
\begin{examinstructions}
Answer 4 questions. You should indicate which 4 problems you wish to be graded. Write your solutions clearly in complete English sentences. You may quote results (within reason) as long as you state them clearly.
\end{examinstructions}
\begin{examproblems}{0}{H2}
\def\ProblemHeadingText{Problem 1.}
\item
This question is about the additive group of rational numbers, denoted by~\(\mathbb{Q}\).
\begin{examsubparts}
\item[\textnormal{(a)}]
Prove that the group~\(\mathbb{Q}\) has no proper subgroups of finite index.
\item[\textnormal{(b)}]
Prove that~\(\mathbb{Q}\) does not have a finite generating set.
\end{examsubparts}
\def\ProblemHeadingText{Problem 2.}
\item
Let~\(G\) be a finite group and~\(N\) a normal subgroup. Show that if~\(p\) is a prime and~\(P\) is a Sylow~\(p\)-subgroup of~\(G\), then \(P\cap N\) is a Sylow~\(p\)-subgroup of~\(N\).
\def\ProblemHeadingText{Problem 3.}
\item
Let~\(G\) be a finite group.
\begin{examsubparts}
\item[\textnormal{(a)}]
Define what is meant by a maximal subgroup of~\(G\).
\item[\textnormal{(b)}]
Let~\(H\) denote the intersection of all the maximal subgroups of~\(G\). Prove that~\(H\) is a characteristic subgroup of~\(G\).
\item[\textnormal{(c)}]
Let~\(H\) be as in (b). Prove that if~\(X\) is a subset of~\(G\) such that \(\langle X\cup H\rangle=G\), then \(\langle X\rangle=G\). (Here \(\langle S\rangle\) denotes the subgroup generated by the set~\(S\).)
\end{examsubparts}
\def\ProblemHeadingText{Problem 4.}
\item
Let~\(G\) be a nonabelian simple group. In \(G\times G\), let \(G_1=\{(g,1)\mid g\in G\}\), \(G_2=\{(1,g)\mid g\in G\}\). Prove that \(G_1\) and \(G_2\) are the only proper nontrivial normal subgroups of \(G\times G\).
\def\ProblemHeadingText{Problem 5.}
\item
Let \(S_n\) and \(A_n\), denote the symmetric group and alternating group of degree~\(n\) respectively, where \(n\geq 2\). Suppose~\(\sigma\) and \(\sigma'\) are elements of \(A_n\) that are in the same \(S_n\)-conjugacy class. Show that they lie in the same \(A_n\)-conjugacy class if and only if there is an odd permutation in the centralizer of~\(\sigma\) in \(S_n\).
\def\ProblemHeadingText{Problem 6.}
\item
This problem has two parts.
\begin{examsubparts}
\item[\textnormal{(a)}]
Compute the order of \(\operatorname{GL}(2,17)\), factored into prime powers.
\item[\textnormal{(b)}]
Prove that a group of order \(7\cdot 17^2\) must be abelian.
\end{examsubparts}
\end{examproblems}
\end{document}
